% Cada entrada é um facto independente da história completa.
\MMRelationshipMapHeading
\begin{itemize}
  \item[] \textit{Origens e os primeiros membros da casa}

  \item \timestamp{Há milhares de anos}, \Character{vampirehead} tornou-se vampira depois de desejar libertar-se da humanidade.
  \item Depois de se tornar vampira, \Character{vampirehead} sentiu-se só e decidiu criar uma família.
  \item \timestamp{Há três mil anos}, \Character{vampirehead} conheceu a humana sedenta de poder \Character{usurper}.
  \item \timestamp{Há três mil anos}, \Character{usurper} concebeu uma poção de sangue de vampiro capaz de transformar um humano num membro da família.
  \item Um vampiro cujo sangue seja usado na poção de transformação pode acrescentar condições ao contrato do novo vampiro.
  \item Durante os mil anos que se seguiram ao seu encontro, \Character{vampirehead} e \Character{usurper} desfrutaram da vida juntas.
  \item Durante esse milénio, a sede de poder de \Character{usurper} começou a desagradar a \Character{vampirehead}.
  \item Em sua casa, \Character{usurper} mantinha notas cifradas, artefactos espalhados e uma desordem deliberada.
  \item \timestamp{Há dois mil anos}, \Character{vampirehead} acolheu \Character{enforcer} na família para fazer cumprir as suas regras.
  \item A autoridade de \Character{enforcer} sobre as regras só era superada pela de \Character{vampirehead}.
  \item \timestamp{Há dois mil anos}, \Character{vampirehead} escreveu as regras da família no \Clue{rulebook}.
  \item Uma votação da família que considere um vampiro culpado de violar uma regra retira-lhe os poderes.
  \item O \Clue{rulebook} torna \Character{usurper} herdeira dos poderes de \Character{vampirehead}, caso \Character{vampirehead} os perca.
  \item \Character{usurper} reparou na possibilidade de herdar os poderes de \Character{vampirehead}.
  \item O \Clue{rulebook} mágico mostrou a \Character{usurper} que \Character{vampirehead} podia criar artefactos poderosos.
  \item \Character{enforcer} escondeu o \Clue{rulebook} algures no castelo.
  \item As regras da família contidas no \Clue{rulebook} eram públicas antes da perda de memória.
  \item Nos encontros de família, \Character{enforcer} repetia algumas regras do \Clue{rulebook}.
  \item \timestamp{Algumas décadas depois da morte de Cristo}, \Character{buttler} juntou-se à família para gerir a casa e as finanças.
  \item \Character{enforcer} dizia abertamente que queria uma mordoma só para si.
  \item \timestamp{No ano 100 d.C.}, \Character{enforcer} apaixonou-se pelo bobo humano \Character{ghoul}.
  \item As regras da família proíbem os vampiros de ter relações com humanos.
  \item \timestamp{No ano 100 d.C.}, \Character{enforcer} pediu a \Character{vampirehead} que transformasse \Character{ghoul} num vampiro.
  \item Só \Character{enforcer} e \Character{ghoul} estiveram presentes na cerimónia secreta de transformação de \Character{ghoul}.
  \item Foi usado o sangue de \Character{enforcer} para transformar \Character{ghoul}.
  \item O contrato de \Character{enforcer} com \Character{ghoul} obriga \Character{ghoul} a obedecer a todas as ordens de \Character{enforcer}.
  \item \Character{enforcer} ordenou a \Character{ghoul} que lhe transmitisse todas as informações relevantes que descobrisse sobre a família.
  \item Na falsa cerimónia pública de transformação de \Character{ghoul}, \Character{buttler} reparou que o frasco dele não continha a poção.
  \item Depois da falsa cerimónia, \Character{vampirehead} pagou a \Character{buttler} para guardar silêncio sobre o frasco.
  \item \Character{usurper} executava por vezes os traidores quando \Character{enforcer} não estava disponível.
  \item Ao executar traidores, \Character{usurper} descobriu que a poção proibida mata vampiros.
  \item O exílio ou a execução de familiares conspiradores ensinou a \Character{usurper} que derrubar \Character{vampirehead} seria difícil.
  \item \timestamp{No ano 500 d.C.}, \Character{supremacist} juntou-se à família.
  \item \Character{supremacist} apaixonou-se secretamente por \Character{vampirehead}.
  \item \Character{supremacist} tratava bem a família por estar apaixonada por \Character{vampirehead}.
  \item \Character{vampirehead} atribuiu a \Character{supremacist} a decisão final sobre a entrada de um humano na família.
  \item \timestamp{No ano 700 d.C.}, \Character{lover} juntou-se à família como protetor.
  \item \timestamp{No ano 1000 d.C.}, \Character{prevlover} juntou-se à família como investigador.
  \item \Character{usurper} deu a \Character{prevlover} o \Clue{clue:truthseeing} para o ajudar a resolver crimes.
  \item \Character{prevlover} e \Character{enforcer} trabalhavam frequentemente juntos em investigações.
  \item \timestamp{No ano 1400 d.C.}, \Character{traitor} juntou-se à família porque ela e \Character{lover} estavam apaixonados.
  \item A natureza de boa samaritana de \Character{traitor} tornou controversa a sua admissão na família.
  \item Depois de entrar na família, \Character{traitor} começou a conspirar com a \outsideentity{igreja}.
  \item \Character{traitor} quer que os vampiros sejam curados ou destruídos e que os humanos vulneráveis sejam protegidos.
  \item Todos conhecem as opiniões de \Character{traitor} contra a família e a consideram hipócrita.
  \item \Character{usurper} brinca publicamente com a ideia de \Character{traitor} voltar a beber a poção proibida.
  \item \Character{usurper} receia que \Character{traitor} impeça a sua sucessão.
  \item Ao longo dos anos, \Character{usurper} ficou conhecida como a familiar prestável.
  \item \Character{supremacist} acredita que os humanos devem ser mortos ou mantidos como gado.
  \item \Character{supremacist} nunca confessou o seu amor a \Character{vampirehead}.
  \item A família trata \Character{ghoul} como uma piada em vez de um contador de piadas.
  \item \Character{ghoul} já não ama \Character{enforcer}.
  \item \Character{ghoul} quer que \Character{enforcer} morra para terminar a obrigação de lhe obedecer.

  \item[] \textit{Os últimos dois anos}

  \item \timestamp{Há dois anos}, \Character{vampirehead} conheceu \Character{newcomer} e convidou-o para uma festa de família.
  \item Nessa festa, \Character{newcomer} conheceu os outros vampiros pela primeira vez.
  \item Todos suspeitam de que \Character{vampirehead} está a cortejar \Character{newcomer}.
  \item Todos esperam que \Character{vampirehead} transforme \Character{newcomer} num membro da família.
  \item As regras da família impedem \Character{vampirehead} de ter uma relação com \Character{newcomer} antes da transformação.
  \item \Character{supremacist} tinha ciúmes de \Character{newcomer}.
  \item \Character{supremacist} opôs-se à entrada de \Character{newcomer} na família.
  \item \Character{vampirehead} ameaçou retirar os poderes a \Character{supremacist}.
  \item Durante os últimos dois anos, \Character{newcomer} e \Character{prevlover} começaram uma relação amorosa secreta.
  \item \timestamp{Há um ano}, \Character{newcomer} começou uma relação amorosa com \Character{lover}.
  \item \timestamp{Há um ano}, \Character{newcomer} terminou a relação com \Character{prevlover}.
  \item Ao deixar \Character{prevlover}, \Character{newcomer} disse que a relação deles estava errada.
  \item \Character{vampirehead} continuou sem saber dos casos amorosos secretos de \Character{newcomer}.
  \item \Character{vampirehead} manteve a distância de \Character{newcomer} para não violar as regras.
  \item \Character{traitor} suspeita de que \Character{lover} conhece os seus planos com a \outsideentity{igreja}.
  \item \Character{traitor} ainda ama \Character{lover}.
  \item Durante o último ano, \Character{ghoul} roubou uma página do \Clue{rulebook}, dando origem à pista \Clue{clue:hierarchy_page}.
  \item \Character{ghoul} roubou a página porque odiava ser lembrado de que tem de obedecer a \Character{enforcer}.
  \item \Character{enforcer} reparou que faltava uma página do \Clue{rulebook} e iniciou uma investigação.
  \item \Character{enforcer} considerou \Character{traitor} a principal suspeita na investigação da página desaparecida.
  \item \Character{enforcer} pediu a \Character{prevlover} que ajudasse a investigar a página desaparecida.
  \item \Character{enforcer} intercetou mensagens invulgares enviadas de e para a casa de \Character{traitor}.
  \item \Character{prevlover} intercetou mensagens invulgares enviadas de e para a casa de \Character{lover}.
  \item As mensagens intercetadas revelam a correspondência de \Character{traitor} com a \outsideentity{igreja}.
  \item As mensagens intercetadas revelam mensagens amorosas entre \Character{lover} e \Character{newcomer}.
  \item \Character{prevlover} guardou as notas sobre as mensagens intercetadas após a investigação.
  \item \Character{enforcer} e \Character{prevlover} não interpretaram as mensagens intercetadas antes da festa.
  \item \timestamp{Dois meses antes da festa}, \Character{vampirehead} marcou \timestamp{16/7/2026} para a festa de transformação de \Character{newcomer}.
  \item \Character{vampirehead} ordenou a todos os vampiros que comparecessem à festa de transformação de \Character{newcomer}.
  \item \timestamp{Umas duas semanas antes da festa}, \Character{usurper} ofereceu-se para ajudar \Character{vampirehead} a prepará-la.
  \item \Character{vampirehead} pediu à família que chegasse duas noites antes para ajudar a preparar a festa.
  \item \Character{usurper} aceitou o convite de \Character{vampirehead} para ajudar a preparar a festa.
  \item \Character{buttler} aceitou o convite de \Character{vampirehead} para ajudar a preparar a festa.
  \item \Character{newcomer} aceitou o convite de \Character{vampirehead} para ajudar a preparar a festa.
  \item Todos os outros vampiros recusaram o convite de \Character{vampirehead} para ajudar a preparar a festa.
  \item Os convites recusados permitiram aos outros vampiros saber quem estaria no castelo antes da festa.

  \item[] \textit{Dois dias antes da festa}

  \item \timestamp{Duas noites antes da festa}, \Character{usurper} chegou ao castelo.
  \item \timestamp{Duas noites antes da festa}, \Character{newcomer} chegou ao castelo.
  \item \timestamp{Duas noites antes da festa}, \Character{vampirehead} já estava no castelo.
  \item \timestamp{Duas noites antes da festa}, \Character{buttler} já estava no castelo.

  \item[] \textit{Na véspera da festa}

  \item \timestamp{À 01:00}, \Character{ghoul} entrou no gabinete de \Character{usurper}.
  \item \timestamp{À 01:00}, \Character{prevlover} entrou no gabinete de \Character{usurper}.
  \item \timestamp{À 01:00}, \Character{traitor} entrou no gabinete de \Character{usurper}.
  \item \timestamp{À 01:00}, \Character{supremacist} entrou no gabinete de \Character{usurper}.
  \item \timestamp{À 01:00}, \Character{ghoul} roubou o ritual ``Mudar o futuro'' do gabinete de \Character{usurper}.
  \item \Character{ghoul} roubou o ritual do futuro porque queria libertar-se de \Character{enforcer}.
  \item \timestamp{À 01:00}, \Character{prevlover} roubou o ritual ``Mudar e conhecer uma pessoa'' do gabinete de \Character{usurper}.
  \item \Character{prevlover} roubou o ritual da memória porque queria recuperar \Character{newcomer}.
  \item \timestamp{À 01:00}, \Character{ghoul} viu \Character{prevlover}.
  \item Ao sair do gabinete, \Character{prevlover} deixou cair a mala que continha o \Clue{clue:truthseeing} e as notas sobre as mensagens intercetadas.
  \item \timestamp{À 01:00}, \Character{traitor} procurou a poção proibida no gabinete de \Character{usurper}, mas não a encontrou.
  \item \timestamp{À 01:00}, \Character{traitor} roubou o ritual ``Fazer com que o que é deles seja de outra pessoa'' do gabinete de \Character{usurper}.
  \item \Character{traitor} acreditava que o ritual roubado podia tirar os bens a outra pessoa.
  \item Na verdade, o ritual roubado desloca todos os objetos que pertencem ao seu alvo.
  \item \timestamp{À 01:00}, \Character{supremacist} roubou a \Clue{clue:truth_potion} do gabinete de \Character{usurper}.
  \item \Character{supremacist} acreditava que a \Clue{clue:truth_potion} mataria humanos, mas não vampiros.
  \item Um vampiro que beba a \Clue{clue:truth_potion} pode obrigar seres conscientes a responder com verdade.
  \item \timestamp{Às 08:00}, \Character{supremacist} colocou metade da \Clue{clue:truth_potion} no vinho reservado para \Character{newcomer}.
  \item O vinho envenenado por \Character{supremacist} tornou-se a pista \Clue{clue:poisoned_wine}.
  \item Depois de envenenar o vinho, \Character{supremacist} guardou a outra metade da \Clue{clue:truth_potion}.
  \item \timestamp{Às 08:00}, \Character{supremacist} disse a \Character{buttler} que tinha mudado de opinião.
  \item A alegada mudança de opinião de \Character{supremacist} dizia respeito a \Character{newcomer}.
  \item \Character{supremacist} falou a \Character{buttler} de um vinho especial.
  \item \Character{supremacist} destinava o vinho especial a \Character{newcomer}.
  \item \Character{supremacist} pagou a \Character{buttler} para manter em segredo a visita e a suposta mudança de opinião.
  \item \timestamp{Às 12:00}, \Character{usurper} substituiu o vinho que \Character{buttler} ia servir por vinho com a poção proibida.
  \item \timestamp{Às 15:00}, \Character{ghoul} realizou ``Mudar o futuro'' e viu a morte de \Character{newcomer}.
  \item O ritual do futuro de \Character{ghoul} deixou os \Clue{clue:futritualleftovers}.
  \item \Character{ghoul} tem de contar a \Character{enforcer} o que descobre sobre a família.
  \item A visão de \Character{ghoul} implica que \Character{newcomer} já tinha sido transformado.
  \item \Character{enforcer} escondeu os roubos de \Character{vampirehead}.
  \item \timestamp{Às 16:00}, \Character{enforcer} telefonou a \Character{prevlover} e ofereceu-se para ignorar o roubo se ele entregasse o ritual.
  \item \Character{enforcer} queria usar o ritual roubado para descobrir por que morreria \Character{newcomer}.
  \item \Character{prevlover} receava que o ritual revelasse a sua antiga relação com \Character{newcomer}.
  \item \Character{prevlover} receava perder os poderes se a relação com \Character{newcomer} fosse revelada.
  \item \timestamp{Às 16:15}, \Character{prevlover} usou o ritual roubado na memória de \Character{enforcer}.
  \item O ritual da memória de \Character{prevlover} deixou os \Clue{clue:memritualleftovers}.
  \item O ritual da memória de \Character{prevlover} apagou memórias aleatórias de todos os vampiros.
  \item O ritual da memória apagou sobretudo recordações do gabinete de \Character{usurper} e dos seus artefactos.
  \item O ritual da memória apagou as recordações da família sobre as regras do \Clue{rulebook}.
  \item \timestamp{Às 17:00}, \Character{traitor} realizou ``Fazer com que o que é deles seja de outra pessoa''.
  \item O ritual de \Character{traitor} espalhou os objetos do gabinete de \Character{usurper} pelo castelo.
  \item O ritual de \Character{traitor} deixou os \Clue{clue:stealritualleftovers}.
  \item \timestamp{Às 21:00}, \Character{lover} telefonou a \Character{newcomer} para que viesse lá fora.
  \item \Character{vampirehead} reparou que \Character{newcomer} saiu do castelo às \timestamp{21:00}.
  \item \Character{traitor} reparou que \Character{lover} não estava em casa às \timestamp{21:00}.
  \item \timestamp{Às 21:00}, \Character{newcomer} disse a \Character{lover} que a relação tinha de acabar.
  \item \Character{newcomer} acreditava que \Character{vampirehead} se oporia ao fim da relação.
  \item Depois da conversa, \Character{newcomer} e \Character{lover} passaram uma última noite juntos.

  \item[] \textit{No dia da festa}

  \item \timestamp{Às 08:00}, \Character{vampirehead} perguntou a \Character{newcomer} onde tinha estado na noite anterior.
  \item \timestamp{Às 08:00}, \Character{newcomer} recusou dizer a \Character{vampirehead} onde tinha estado.
  \item \timestamp{Às 11:00}, \Character{lover} e \Character{newcomer} conversaram lá fora.
  \item \timestamp{Às 11:00}, \Character{lover} disse a \Character{newcomer} que não queria que a relação terminasse.
  \item \timestamp{Às 11:00}, \Character{buttler} viu a conversa no exterior.
  \item \timestamp{Às 11:00}, \Character{usurper} viu a conversa no exterior.
  \item \Character{buttler} e \Character{usurper} não se aperceberam uma da outra durante a conversa no exterior.
  \item Depois da conversa no exterior, \Character{buttler} ofereceu a \Character{lover} o seu silêncio em troca de um pagamento.
  \item \Character{lover} pagou a \Character{buttler} pelo silêncio sobre a conversa no exterior.
\end{itemize}
